% ===========================================================================
%  Document macros shared by all article-style documents.
%
%  Loaded last by preamble.tex, so every package and the \ifnpprint
%  switch are already available here.
%
%  The macros you will actually use while writing content:
%
%      \defi{...}     green  "Erklärung"  box   -- a definition / explanation
%      \methode{...}  blue   "Vorgehen"   box   -- a recipe, step by step
%      \bsp{...}      orange "Beispiel"   box   -- a worked example
%      \wichtig{...}  red    "Wichtig"    box   -- something not to forget
%      \vertief{...}  brown  "Vertiefung" box   -- optional deeper background
%
%      \pic{f}        \picS{f}{width}     \picC{f}{width}{caption}
%                     include gemeinsam/bilder/f, relative to \linewidth
%
%      \n \sbreak \lbreak     vertical space,  \fat{...}  bold,  \qu{...} quotes
%
%  In the print variant (\npstyle = print) every box is black and \red / \green
%  / ... fall back to bold text, so nothing depends on colour being visible.
% ===========================================================================

\newcommand{\qu}[1]{\enquote{#1}}

% ---------------------------------------------------------------------------
% Where the images live.  Documents that sit deeper in the tree only have to
% set \npbase; the picture macros follow automatically.
% ---------------------------------------------------------------------------
\providecommand{\nppicdir}{\npbase gemeinsam/bilder/}

% ===========================================================================
% Info boxes
% ===========================================================================

\definecolor{lightgray}{gray}{0.95}
\definecolor{middlegray}{gray}{0.75}
\definecolor{darkgreen}{rgb}{0, 0.5, 0}

% A vertical coloured bar down the left margin of a block of text.
\newenvironment{MyLeftBar}[1]{%
	\def\FrameCommand##1{\textcolor{#1}{\vrule width 2pt} \hspace{9pt}##1}%
	\MakeFramed {\advance\hsize-\width \FrameRestore}}%
{\endMakeFramed}

\newenvironment{LeftBar}[1]{%
	\def\FrameCommand##1{\textcolor{#1}{\vrule width 2pt} \hspace{9pt}##1}%
	\MakeFramed {\advance\hsize-\width \FrameRestore}}%
{\endMakeFramed}

\newenvironment{InfoBox}[3]{
	\begin{LeftBar}{#2}\hspace{-25pt}\colorbox{#2}{\textbf{\textcolor{white}{#1}}}{\indent{}#3\vspace{0.5em}\\noindent}}{\end{LeftBar}}

% notebox{colour}{label}{body} -- the coloured bar with a title tag.
% NOTE: this used to be called \note, which clashes with beamer's \note and
% broke every presentation that loaded this file.
\newcommand*{\notebox}[3]{%
	\begin{MyLeftBar}{#1}\begin{samepage}%
		\hspace{-25pt}\colorbox{#1}{\textbf{\textcolor{white}{#2}}}{\indent #3}%
	\end{samepage}\end{MyLeftBar}}

% The five box flavours.  \npboxcolor picks the accent colour, or black when
% the document is built for grayscale printing.
\ifnpprint
	\newcommand*{\npboxcolor}[1]{black}
\else
	\newcommand*{\npboxcolor}[1]{#1}
\fi

\newcommand*{\methode}[1]{\notebox{\npboxcolor{blue}}{Vorgehen}{\hspace{1em} #1}}
\newcommand*{\bsp}[1]{\notebox{\npboxcolor{orange}}{Beispiel}{\hspace{1em} #1}}
\newcommand*{\wichtig}[1]{\notebox{\npboxcolor{red}}{Wichtig}{\hspace{1em} #1}}
\newcommand*{\defi}[1]{\notebox{\npboxcolor{darkgreen}}{Erklärung}{\hspace{1em} #1}}
\newcommand*{\vertief}[1]{\notebox{\npboxcolor{brown}}{Vertiefung}{\hspace{1em} #1}}

% ===========================================================================


% =========================================================
%\green{Defintion}
%\red{Wichtig}
%\orange{Beispiel(e)}
%\blue{Erklärung/Methodik und im Anderen Aufgaben}

% =========================================================
% Ref Commands
% =========================================================

% References:
\newcommand{\refFigure}[1]{\hyperref[#1]{Figur \ref{#1}}}
\newcommand{\refTable}[1]{\hyperref[#1]{Tabelle \ref{#1}}}
\newcommand{\refChapter}[1]{\hyperref[#1]{Kapitel \ref{#1}}}
\newcommand{\refEquation}[1]{\hyperref[#1]{Gleichung \ref{#1}}}

% multiple refs:
\newcommand{\refChapters}[2]{\hyperref[#1]{Kapitel \ref{#1}} \hyperref[#2]{und \ref{#2}}}


% =========================================================
% Text Commands
% =========================================================

\newcommand{\lingform}[1]{{\it{#1}}}		% linguistic form
\newcommand{\scare}[1]{`#1'} 
\newcommand{\code}[1]{{\tt{#1}}}

\newcommand{\marK}[2]{\fcolorbox{#1}{white}{#2}}

% Inline highlighting.  In the print variant colour carries no information,
% so every one of these degrades to bold text instead of a grey smudge.
\ifnpprint
	\def\red#1{\textbf{#1}}
	\def\blue#1{\textbf{#1}}
	\def\orange#1{\textbf{#1}}
	\def\green#1{\textbf{#1}}
	\def\gray#1{#1}
	\def\lightgray#1{#1}
	\def\black#1{#1}
\else
	\def\red#1{\textcolor{red!90!black}{#1}}
	\def\blue#1{\textcolor{blue!75!black}{#1}}
	\def\orange#1{\textcolor{orange!95!black}{#1}}
	\def\green#1{\textcolor{green!50!black}{#1}}
	\def\gray#1{\textcolor{gray!90!black}{#1}}
	\def\lightgray#1{\textcolor{lightgray!90!black}{#1}}
	\def\black#1{\textcolor{black}{#1}}
\fi

\renewcommand{\Re}{\mathbb{R}}

\renewcommand*{\thefootnote}{*}

\def\checkc{\text{\textcolor{green!50!black}{\ding{51}}}}
\def\checkx{\text{\textcolor{red}{\ding{55}}}}

\newcommand{\fat}[1]{\textbf{#1}}
\newcommand{\ftitle}[2]{\huge #2{#1}\vspace{.5in}\normalsize}

\newcommand{\sbreak}{\vspace{0.1in}}
\newcommand{\lbreak}{\vspace{0.3in}}

\newcommand{\noted}{\sbreak \\}
\newcommand{\n}{\sbreak \\}

\newcommand{\itemf}{\vfill \item}

\newcommand{\task}[1]{\newpage \section{#1}}
\newcommand{\sol}[1]{\green{#1}}

% Add another entry
\newcommand{\aae}{; \hspace{0.25cm} }

% Standard picture
\newcommand{\pic}[1]{
	\begin{figure}[htbp]
		\centering
		\includegraphics[width=0.45\linewidth]{\nppicdir #1}
	\end{figure}
}
\newcommand{\picS}[2]{
	\begin{figure}[htbp]
		\centering
		\includegraphics[width=#2\linewidth]{\nppicdir #1}
	\end{figure}
}
\newcommand{\picC}[3]{
	\begin{figure}[htbp]
		\centering
		\includegraphics[width=#2\linewidth]{\nppicdir #1}
		\caption{#3}
	\end{figure}
}

% =========================================================
% Table commands:
% =========================================================

\newcommand{\greyMidRule}{\arrayrulecolor{black!30}\midrule}

% =========================================================
% Specific math variables:
% =========================================================

\newcommand{\fabs}{f_{\text{abs}}}
\newcommand{\frel}{f_{\text{rel}}}

\newcommand{\maxError}{\epsilon_{\text{Ma}}}

\newcommand{\mddots}{\reflectbox{$\ddots$}}


\newcommand{\umax}{_\text{max}}
\newcommand{\umin}{_\text{min}}

% =========================================================
% Math Commands
% =========================================================

% Should be renamed or removed
\newcommand{\st}[1]{^{(#1)}}

% vertical line in the middle:
\newcommand{\separator}[1]{\; #1 \;}
\newcommand{\sep}{\separator{|}}
\newcommand{\sepArrow}{\separator{\Longleftrightarrow}}

% BigO
\newcommand{\bigO}{\mathcal{O}} % Big-O notation

% Norm with double |

% Matrix Transpose and inverted
\newcommand{\trans}{^{\mathrm{\mathsmaller T}}}
\newcommand{\inv}{^{\mathrm{\mathsmaller{\text{-}1}}}}
\newcommand{\unitM}{\mathbb{1}}


% informal descriptions:
\newcommand{\uleftDown}{_\text{linksUnten}}
\newcommand{\uleftUp}{_\text{linksOben}}
\newcommand{\uleft}{_\text{links}}
\newcommand{\uleftDiagonal}{_\text{ganzLinksUnten}}
\newcommand{\uleftUpDiagonal}{_\text{ganzLinksOben}}

% Soll sein
\newcommand{\shouldEq}{ \stackrel{!}{=} }

% Formelnamen:
\newcommand{\QT}{Q_{\text{T}}}
\newcommand{\QTS}{Q_{\text{TS}}}
\newcommand{\QS}{Q_{\text{S}}}
\newcommand{\QSS}{Q_{\text{SS}}}
\newcommand{\RTS}{R_{\text{TS}}}
\newcommand{\RSS}{R_{\text{SS}}}
%--------------------------

% correct strange spacing around delimiters
\renewcommand{\l}{\mathopen{}\mathclose\bgroup\left}
\renewcommand{\r}{\aftergroup\egroup\right}

% norms
\newcommand{\norm}[1]{\l\lVert #1 \r\rVert}
\newcommand{\pnorm}[2]{\norm{#1}_{#2}}

% column vectors		(z.B. \colvec{3}{x_1}{\vdots}{x_n})
\newcount\colveccount
\newcommand*\colvec[1]{
	\global\colveccount#1
	\begin{pmatrix}
		\colvecnext
	}
	\def\colvecnext#1{
		#1
		\global\advance\colveccount-1
		\ifnum\colveccount>0
		\\
		\expandafter\colvecnext
		\else
	\end{pmatrix}
	\fi
}

\newcommand{\nvector}[1]{
	\begin{pmatrix}
		#1
	\end{pmatrix}
}

\newcommand{\nmatrix}[1]{
	\begin{bmatrix}
		#1
	\end{bmatrix}
}

% congruence modulo n
\def\congr#1{\equiv_{#1}}
\def\ncongr#1{\not \equiv_{#1}}

\def\dann#1#2{
	\begin{array}{lll}
		#1 & \Longrightarrow & #2
	\end{array}
}

\def\gdw#1#2{
	\begin{array}{lll}
		#1 & \Longleftrightarrow & #2
	\end{array}
}

\def\defgdw#1#2{
	\begin{array}{lll}
		#1 & :\Longleftrightarrow & #2
	\end{array}
}

\def\nLongrightarrow{\centernot \Longrightarrow}
\def\nLongleftrightarrow{\centernot \Longleftrightarrow}
\def\nforall{\centernot \forall}
\def\nexists{\centernot \exists}
\def\nsqsubseteq{\centernot \sqsubseteq}
\def\nsubset{\centernot \subset}
\def\nequiv{\centernot \equiv}

\def\symdif{\triangle}
\def\imp{\rightarrow}
\def\bik{\leftrightarrow}
\def\xor{\oplus}
\def\nand{\, \bar{\land}\,}
\def\nor{\,\bar{\lor}\,}


\def\iv#1{\overset{\text{I.V.}}{#1}}

% equivalence relations
\newcommand{\quotientset}[2]{{{#1}/{#2}}}			% curly brackets around #1 and #2 because of some spacing problems
\renewcommand{\index}[2]{\card{\quotientset{#1}{#2}}}

% discrete intervals (with or without zero)
\newcommand{\di}[2]{\set{#1,\ldots,#2}}				% {a,...,b}
\newcommand{\dizero}[1]{\sb{#1}_0}					% [n]_0 = {0,...,n}
\newcommand{\dione}[1]{\sb{#1}}						% [n] = {1,...,n}

% cardinality
\newcommand{\card}[1]{\l|#1\r|}
\newcommand{\cardalt}[2]{
	\ifthenelse{#1 = 1}{\operatorname{card}\(#2\)}{
		\ifthenelse{#1 = 2}{\#\(#2\)}{
			\ifthenelse{#1 = 3}{n\(#2\)}{
				\ifthenelse{#1 = 4}{\overline{\overline{#2}}}{
					\text{TODO: Command not defined!}
				}
			}
		}
	}
}

% matrix/vector entries
\newcommand{\mentry}[3]{#1_{#2 #3}}

% word length
\newcommand{\wordlength}[1]{\l|#1\r|}
\newcommand{\length}[1]{\l|#1\r|}

% number of characters in word
\newcommand{\numchar}[2]{\l|#2\r|_{#1}}
\newcommand{\numcharalt}[3]{
	\ifthenelse{#1 = 1}{\#_{#2}\rb{#3}}{
		\text{TODO: Command not defined!}
	}
}

% complement set
\newcommand{\compl}[1]{\overline{#1}}

% power set
\newcommand{\powersetempty}[0]{\mathcal{P}}
\newcommand{\powerset}[1]{\powersetempty\l(#1\r)}
\newcommand{\powersetalt}[2]{
	\ifthenelse{#1 = 1}{P\(#2\)}{
		\ifthenelse{#1 = 2}{\mathbb{P}\(#2\)}{
			\ifthenelse{#1 = 3}{\wp\(#2\)}{
				\ifthenelse{#1 = 4}{2^{#2}}{
					\text{TODO: Command not defined!}
				}
			}
		}
	}
}

% sets
\newcommand{\set}[1]{\left\{#1\right\}}
\newcommand{\suchthat}[0]{\, \middle| \,}
\newcommand{\suchthatalt}[1]{
	\ifthenelse{#1 = 1}{:}{
		\ifthenelse{#1 = 2}{\, ;}{
			\text{TODO: Command not defined!}
		}
	}
}

% number sets
\newcommand{\booleans}[0]{\mathbb{B}}			% Perhaps call it \binaries and defined it as \residuals{2} instead.
\newcommand{\naturals}[0]{\mathbb{N}}
\newcommand{\wholes}[0]{\mathbb{N}_0}
\newcommand{\integers}[0]{\mathbb{Z}}
\newcommand{\rationals}[0]{\mathbb{Q}}
\newcommand{\reals}[0]{\mathbb{R}}
\newcommand{\complexes}[0]{\mathbb{C}}
\newcommand{\residuals}[1]{\mathbb{Z}_{#1}}
\newcommand{\primes}[0]{\mathbb{P}}

% multisets
\newcommand{\multiset}[1]{\left\{\!\left\vert #1 \right\vert\!\right\}}
\newcommand{\multisetalt}[2]{
	\ifthenelse{#1 = 1}{\set{#2}}{
		\text{TODO: Command not defined!}
	}
}

% concatenation
\newcommand{\concat}[0]{\circ}
\newcommand{\concatalt}[2]{
	\ifthenelse{#1 = 1}{}{
		\ifthenelse{#1 = 2}{\parallel}{
			\text{TODO: Command not defined!}
		}
	}
}

% continous intervals (o = open, c = closed, i = infinity)
\newcommand{\cc}[2]{\l[#1,#2\r]}
\newcommand{\co}[2]{\l[#1,#2\r)}
\newcommand{\ci}[1]{\l[#1,\infty\r]}
\newcommand{\oc}[2]{\l(#1,#2\r]}
\newcommand{\oo}[2]{\l(#1,#2\r)}
\newcommand{\oi}[1]{\l(#1,\infty\r)}
\newcommand{\ic}[1]{\l(-\infty,#1\r]}
\newcommand{\io}[1]{\l(-\infty,#1\r]}
\newcommand{\ii}[0]{\l(-\infty,\infty\r)}

% Fourier
\newcommand{\fourier}[3]{\mathcal{F}_{#1,#2}\l(#3\r)}
\newcommand{\Pveca}[1]{P_{\vec{a}}\l(#1\r)}

% falling and rising factorials
\newcommand{\fallfac}[2]{#1^{\underline{#2}}}
\newcommand{\risefac}[2]{#1^{\overline{#2}}}

% binomial coefficient
\newcommand{\bin}[2]{\binom{#1}{#2}}

% stirling numbers
\newcommand{\stirlf}[2]{\genfrac{[}{]}{0pt}{}{#1}{#2}}
\newcommand{\stirls}[2]{\genfrac{\{}{\}}{0pt}{}{#1}{#2}}

% sequences
\newcommand{\seq}[3]{
	\ifthenelse{#3 = 0}{(#1_#2)_{#2 \in \wholes}}{
		\ifthenelse{#3 = 1}{(#1_#2)_{#2 \in \naturals}}{
			(#1_#2)_{#2 \geq #3}
		}
	}
}

% series
\newcommand{\ser}[3]{
	\sum_{#1=#2}^{\infty} #3
}

% general series
\newcommand{\genser}[1]{
	\ser{k}{1}{#1_k}
}

% zero set
\newcommand{\zeroset}[1]{\preimage{#1}{0}}		% set of all roots of a function #1

% preimage
\newcommand{\preimage}[2]{#1^{-1}\l(#2\r)}		% set of all x suchthat #1(x) = #2

% brackets (round, square, curly, angled)
\renewcommand{\(}[0]{\l(}						% \(...\) was defined for $...$ as \[...\] for $$...$$
\renewcommand{\)}[0]{\r)}
\newcommand{\rb}[1]{\l(#1\r)}
\renewcommand{\sb}[1]{\l[#1\r]}
\newcommand{\cb}[1]{\l\{#1\r\}}
\newcommand{\ab}[1]{\l\langle#1\r\rangle}

% generating set
\newcommand{\generatingset}[1]{\ab{#1}}

% floor and ceiling
\newcommand{\ceil}[1]{\l\lceil#1\r\rceil}
\newcommand{\floor}[1]{\l\lfloor#1\r\rfloor}


% absolute value
\newcommand{\abs}[1]{\l|#1\r|}

% integrals
\newcommand{\integral}[4]{\int_{#1}^{#2} #3 \ \mathrm{d}#4}
\newcommand{\antiderivative}[2]{\integral{}{}{#1}{#2}}
\renewcommand{\d}[1]{\ \mathrm{d}#1}

% logarithms
\newcommand{\ld}[0]{\operatorname{ld}}
\newcommand{\lb}[0]{\operatorname{lb}}

% signum function
\newcommand{\sgn}[0]{\operatorname{sgn}}

% fixed set
\newcommand{\fix}[0]{\operatorname{fix}}		% fix(f) = {x|f(x) = x}

% closures
\newcommand{\refl}[0]{\operatorname{ref}}
\newcommand{\symm}[0]{\operatorname{sym}}

% sets of functions
\newcommand{\inj}[0]{\operatorname{inj}}
\newcommand{\surj}[0]{\operatorname{surj}}
\newcommand{\bij}[0]{\operatorname{bij}}

% fields
\newcommand{\charac}[0]{\operatorname{char}}
\newcommand{\GF}[0]{\operatorname{GF}}

% least common multiple and greatest common divisor
\newcommand{\ggT}[0]{\operatorname{ggT}}
\newcommand{\kgV}[0]{\operatorname{kgV}}

% average
\newcommand{\avg}[0]{\operatorname{avg}}

% order
\newcommand{\ord}[0]{\operatorname{ord}}

% ----------
% round 
\newcommand{\rd}[0]{\operatorname{rd}}
\newcommand{\rdM}[0]{\operatorname{rd_M}}

% condition value
\newcommand{\cond}[0]{\operatorname{cond}}

% complex numbers
\newcommand{\realPart}{\operatorname{Re}}
\newcommand{\imaginaryPart}{\operatorname{Im}}

% Fourier transformation
\newcommand{\DFT}{\operatorname{DFT}}
\newcommand{\IDFT}{\operatorname{IDFT}}
\newcommand{\FFT}{\operatorname{FFT}}
\newcommand{\BFO}{\operatorname{BFO}}

% diagonal entries of a matrix
\newcommand{\diag}{\operatorname{diag}}
% --------

% sequent calculus
\newcommand{\FL}[0]{\operatorname{FL}}
\newcommand{\Nat}[0]{\operatorname{Nat}}

% true and false
\newcommand{\false}[0]{\operatorname{false}}
\newcommand{\true}[0]{\operatorname{true}}

% probability function (alone, with argument, with index, conditional)

\newcommand{\Pempty}[0]{\operatorname{Pr}}		% avoid \Pr, since \Pr_A looks bad in display math mode!
\renewcommand{\P}[1]{\Pempty\sb{#1}}
\newcommand{\Palt}[2]{
	\ifthenelse{#1 = 1}{P\rb{#1}}{
		\ifthenelse{#1 = 2}{\mathbb{P}\rb{#1}}{
			\text{TODO: Command not defined!}
		}
	}
}
\newcommand{\Pindex}[2]{\Pempty_{#1}\sb{#2}}
\newcommand{\Pcond}[2]{\P{#1 \middle| #2}}
\newcommand{\Pcondalt}[3]{
	\ifthenelse{#1 = 1}{\Pindex{#3}{#2}}{
		\text{TODO: Command not defined!}
	}
}

% probability density function
\newcommand{\pdf}[1]{f_{#1}}
\newcommand{\pdfcond}[3]{\pdf{\left.#1 \middle| #2\right.}\rb{#3}}		% f_{#1|#2}(#3)
\newcommand{\pdfcondalt}[4]{
	\ifthenelse{#1 = 1}{\pdf{#2}\rb{\left.#4 \middle| #3\right.}}{
		\text{TODO: Command not defined!}
	}
}

% cumulative distribution function
\newcommand{\cdf}[1]{F_{#1}}
\newcommand{\cdfcond}[3]{\cdf{\left.#1 \middle| #2\right.}\rb{#3}}		% F _{#1|#2}(#3)
\newcommand{\cdfcondalt}[4]{
	\ifthenelse{#1 = 1}{\cdf{#2}\rb{\left.#4 \middle| #3\right.}}{
		\text{TODO: Command not defined!}
	}
}

% expected value
\newcommand{\Eempty}[0]{\mathbb{E}}
\newcommand{\E}[1]{\Eempty\sb{#1}}
\newcommand{\Eindex}[2]{\Eempty_{#1}\sb{#2}}
\newcommand{\Econd}[2]{\Eempty\sb{#1 \middle| #2}}	% E[#1|#2]

% variance
\newcommand{\Varempty}[0]{\operatorname{Var}}
\newcommand{\Var}[1]{\Varempty\sb{#1}}
\newcommand{\Varindex}[2]{\Varempty_{#1}\sb{#2}}
\newcommand{\Varcond}[2]{\Varempty\sb{#1 \middle| #2}}	% Var[#1|#2]

\newcommand{\MSE}[0]{\operatorname{MSE}}
\newcommand{\Bias}[0]{\operatorname{Bias}}

% probability distributions
\newcommand{\Ber}[0]{\operatorname{Ber}}
\newcommand{\Bin}[0]{\operatorname{Bin}}
\newcommand{\Poi}[0]{\operatorname{Poi}}
\newcommand{\Geo}[0]{\operatorname{Geo}}
\newcommand{\NegBin}[0]{\operatorname{NegBin}}
\newcommand{\Hyp}[0]{\operatorname{Hyp}}
\newcommand{\NegHyp}[0]{\operatorname{NegHyp}}
\newcommand{\Nor}[0]{\operatorname{\mathcal{N}}}
\newcommand{\Exp}[0]{\operatorname{Exp}}
\newcommand{\Erl}[0]{\operatorname{Erl}}
\newcommand{\ConUni}[0]{\operatorname{Uni}}
\newcommand{\DisUni}[0]{\operatorname{Uni}}
\newcommand{\Uni}[0]{\operatorname{Uni}}					% TODO: cambiar por \UniDis y \UniCon

% relations
\newcommand{\id}[1]{\operatorname{id}_{#1}}					% Identitätsrelation
\newcommand{\all}[1]{#1 \times #1}							% Allrelation

% =========================================================
% TikZ Commands
% =========================================================

\newcommand{\tikzPoint}[3]{
	\node[label={#1:{#2}},circle,fill,inner sep=2pt] at (axis cs:#3) {};
}
\newcommand{\tikzPointLeft}[2]{
	\tikzPoint{180}{#1}{#2}
}
\newcommand{\tikzPointRight}[2]{
	\tikzPoint{0}{#1}{#2}
}
\newcommand{\tikzPointUp}[2]{
	\tikzPoint{90}{#1}{#2}
}
\newcommand{\tikzPointDown}[2]{
	\tikzPoint{270}{#1}{#2}
}
\newcommand{\tikzPointUL}[2]{
	\tikzPoint{135}{#1}{#2}
}
\newcommand{\tikzPointUR}[2]{
	\tikzPoint{45}{#1}{#2}
}
\newcommand{\tikzPointDL}[2]{
	\tikzPoint{225}{#1}{#2}
}
\newcommand{\tikzPointDR}[2]{
	\tikzPoint{315}{#1}{#2}
}

%[red, very thick, |-] (0,0) -- (22,0);
\def \tikzPattern{on 1pt off 3pt on 3pt off 3pt}%dash pattern=on 1pt off 3pt on 3pt off 3pt}

\newcommand{\tikzLineDotted}[3]{
	\draw[gray, dotted] (axis cs:#2) -- node[left]{#1} (axis cs:#3);
}
\newcommand{\tikzLine}[4]{
	\draw[#1] (axis cs:#3) -- node[left]{#2} (axis cs:#4);
}
\newcommand{\tikzLineUp}[4]{
	\draw[#1] (axis cs:#3) -- node[above]{#2} (axis cs:#4);
}
\newcommand{\tikzNormalLine}[3]{
	\tikzLine{black, thick}{#1}{#2}{#3}
}
\newcommand{\tikzFatLine}[3]{
	\tikzLine{blue, very thick}{#1}{#2}{#3}
}
\newcommand{\tikzFatterLine}[3]{
	\tikzLine{blue, line width=0.65mm}{#1}{#2}{#3}
}
\newcommand{\tikzMegaFatLine}[3]{
	\tikzLine{blue, line width=0.95mm}{#1}{#2}{#3}
}

%	% #1 = position x
%	% #2 = position y
%	%\def \scaler{0.5}%#4
%	}%
%	%\tikzLine{blue}{}{(#1 - 0.5), (#2 - (0.5*#3)}{(#1 + 0.5), (#2 + (0.5*#3)}%
%}%

\newcommand{\tikzArrow}[3]{
	\draw[very thick, ->] (axis cs:#2) -- node[above]{#1} (axis cs:#3);
}
% Angles:
% 0 -> right, 90 = up, 270 = down, 180
\newcommand{\tikzCircle}[3]{
	\draw(axis cs:#2) circle[#1, radius=#3];
}

\newcommand{\tikzTrigonometric}{
	xtick={-3.14159, -1.5708, 0,1.5708,3.14159,4.7123889},
	xticklabels={$-\pi$, $-\frac{\pi}{2}$, $0$,$\frac{\pi}{2}$,$\pi$,$\frac{3\pi}{2}$},
}


\newcommand{\tikzFigure}[2]{
	\begin{figure}[ht]
		\centering
		\begin{tikzpicture}
			#1
		\end{tikzpicture}
		\caption{#2}
	\end{figure}
}
\newcommand{\tikzFigureStandard}[4]{
	\tikzFigure{
		\begin{axis}[
			axis lines=left,
			xlabel=$x$,
			ylabel=$y$,
			enlargelimits,
			%xtick={-3.14159, -1.5708, 0,1.5708,3.14159,4.7123889},
			%xticklabels={$-\pi$, $-\frac{\pi}{2}$, $0$,$\frac{\pi}{2}$,$\pi$,$\frac{3\pi}{2}$},
			%ymax=2,
			%ymin=-0.5]
			%[domain=0:1,legend pos=outer north east]
			\addplot[domain=#1:#2] {#3} node[left]{};
%				( 0, 4 / pi )
%				( pi, 0 )
%				( 0, 0)
%				( 0, 4 / pi)
%			};
		\end{axis}
	}{#4}
}

\newcommand{\annotation}[4]{	% color, position, annotation, word
	\tikz[remember picture,baseline=(word.base)]{
		\node [inner sep=0cm, outer sep=0.1cm, text depth=.25ex, text height=1.5ex] (word) {#4};
	}
	\tikz[overlay,remember picture]{
		\footnotesize
		\def\scalefactor{0.625};
		\IfEqCase{#2}{
			{n}{\path (word)+(0,\scalefactor) node [above,align=center,#1] (annotation) {#3};
				\draw [->,align=center,#1] (annotation.south) to (word);
			}
			{ne}{\path (word)+(\scalefactor,\scalefactor) node [above right,align=center,#1] (annotation) {#3};
				\draw [->,align=center,#1] (annotation.south west) to [bend right=30](word);
			}
			{e}{\path (word)+(\scalefactor,0) node [right,align=center,#1] (annotation) {#3};
				\draw [->,align=center,#1] (annotation.west) to (word);
			}
			{se}{\path (word)+(\scalefactor,-\scalefactor) node [below right,align=center,#1] (annotation) {#3};
				\draw [->,align=center,#1] (annotation.north west) to [bend left=30] (word);
			}
			{s}{\path (word)+(0,-\scalefactor) node [below,align=center,#1] (annotation) {#3};
				\draw [->,align=center,#1] (annotation.north) to (word);
			}
			{sw}{\path (word)+(-\scalefactor,-\scalefactor) node [below left,align=center,#1] (annotation) {#3};
				\draw [->,align=center,#1] (annotation.north east) to [bend right=30] (word);
			}
			{w}{\path (word)+(-\scalefactor,0) node [left,align=center,#1] (annotation) {#3};
				\draw [->,align=center,#1] (annotation.east) to (word);
			}
			{nw}{\path (word)+(-\scalefactor,\scalefactor) node [above left,align=center,#1] (annotation) {#3};
				\draw [->,align=center,#1] (annotation.south east) to [bend left=30] (word);
			}
		}[\PackageError{annotation}{Undefined option to annotation: #3}{Try integer from the set [12].}]
	}
}

% convex paths in TikZ
\newcommand{\convexpath}[2]{
	[   
	create hullcoords/.code={
		\global\edef\namelist{#1}
		\foreach [count=\counter] \nodename in \namelist {
			\global\edef\numberofnodes{\counter}
			\coordinate (hullcoord\counter) at (\nodename);
		}
		\coordinate (hullcoord0) at (hullcoord\numberofnodes);
		\pgfmathtruncatemacro\lastnumber{\numberofnodes+1}
		\coordinate (hullcoord\lastnumber) at (hullcoord1);
	},
	create hullcoords
	]
	($(hullcoord1)!#2!-90:(hullcoord0)$)
	\foreach [
	evaluate=\currentnode as \previousnode using \currentnode-1,
	evaluate=\currentnode as \nextnode using \currentnode+1
	] \currentnode in {1,...,\numberofnodes} {
		let \p1 = ($(hullcoord\currentnode) - (hullcoord\previousnode)$),
		\n1 = {atan2(\x1,\y1) + 90},
		\p2 = ($(hullcoord\nextnode) - (hullcoord\currentnode)$),
		\n2 = {atan2(\x2,\y2) + 90},
		\n{delta} = {Mod(\n2-\n1,360) - 360}
		in 
		{arc [start angle=\n1, delta angle=\n{delta}, radius=#2]}
		-- ($(hullcoord\nextnode)!#2!-90:(hullcoord\currentnode)$) 
	}
}